\begin{thebibliography}{99}
\bibitem{WalkeV1} R. M. Walke,
\emph{Twisted moments of the Legendre family, rigid Calabi--Yau modular forms,
and a weight-five motive of conductor 1296}, Zenodo preprint v1,
28 October 2026. \url{https://doi.org/10.5281/zenodo.23719224}.
This is the historical public predecessor.
\bibitem{SGA1} A. Grothendieck et al.,
\emph{Revetements etales et groupe fondamental (SGA1)}, Expose XIII,
Definition 2.1.1 and Corollary 2.9, recomposed edition.
\url{https://arxiv.org/pdf/math/0206203}.
\bibitem{SGA4} M. Artin, A. Grothendieck and J.-L. Verdier (eds.),
\emph{Theorie des topos et cohomologie etale des schemas (SGA4)}, tome 3,
Expose XVIII, Theorem 2.9, Var. 2, and Theorem 3.2.5.
\url{https://www.cmls.polytechnique.fr/perso/laszlo/sga4/SGA4-3/sga43.pdf}.
\bibitem{Katz} N. M. Katz,
\emph{Convolution and Equidistribution: Sato--Tate Theorems for
Finite-Field Mellin Transforms}, Chapter 15.
\url{https://web.math.princeton.edu/~nmk/mellin331.pdf}.
\bibitem{Tate} J. Tate, \emph{A review of non-Archimedean elliptic functions},
section on points of finite order, pp. 18--19 of the manuscript.
\url{https://people.math.harvard.edu/~ctm/home/text/others/tate/tate-unif/tate-unif.pdf}.
\bibitem{Milne} J. S. Milne, \emph{Lectures on Etale Cohomology},
especially Theorem 24.1 and the sections on adic cohomology.
\url{https://www.jmilne.org/math/CourseNotes/LECc.pdf}.
\bibitem{StacksTrace} The Stacks Project, Theorem 64.20.2, tag 03V0,
finite-field adic trace formula. \url{https://stacks.math.columbia.edu/tag/03V0}.
\bibitem{StacksFinite} The Stacks Project, Proposition 59.55.2, tag 03QP,
vanishing of finite higher direct images.
\url{https://stacks.math.columbia.edu/tag/03QP}.
\bibitem{StacksDescent} The Stacks Project, Section 59.104, tag 0GEX,
descent for etale sheaves. \url{https://stacks.math.columbia.edu/tag/0GEX}.
\bibitem{DGP} L. Dieulefait, L. Guerberoff and A. Pacetti,
\emph{Proving modularity for a given elliptic curve over an imaginary
quadratic field}, Theorem 4.1. \url{https://arxiv.org/abs/0804.2302}.
\bibitem{Sturm} J. Sturm, \emph{On the congruence of modular forms},
Lecture Notes in Math. \textbf{1240}, Springer, 1987, 275--280,
Theorem~1. \url{https://doi.org/10.1007/BFb0072985}.
\bibitem{LTYZ} L. Long, F.-T. Tu, N. Yui and W. Zudilin,
\emph{Supercongruences for rigid hypergeometric Calabi--Yau threefolds},
Adv. Math. \textbf{393} (2021), 108058.
\url{https://arxiv.org/abs/1705.01663}.
\bibitem{KS} M. Kazalicki and A. J. Scholl,
\emph{Modular forms, de Rham cohomology and congruences},
Trans. Amer. Math. Soc. \textbf{368} (2016), 7097--7117, Sections 7--8.
\url{https://www.dpmms.cam.ac.uk/~ajs1005/preprints/mfdrc.pdf}.
\bibitem{Scholl} A. J. Scholl,
\emph{On some $\ell$-adic representations of
$\operatorname{Gal}(\overline\Q/\Q)$ attached to noncongruence subgroups},
Bull. London Math. Soc. \textbf{38} (2006), 561--567.
\url{https://www.dpmms.cam.ac.uk/~ajs1005/preprints/galois-2004.pdf}.
\bibitem{ALLL} A. O. L. Atkin, W.-C. W. Li, T. Liu and L. Long,
\emph{Galois representations with quaternion multiplication associated
to noncongruence modular forms}, Trans. Amer. Math. Soc.
\textbf{365} (2013), 6217--6242.
\url{https://www.math.purdue.edu/~tongliu/pub/tran6019.pdf}.
\bibitem{Grove} B. Grove, \emph{Hypergeometric moments and Hecke trace formulas},
Ramanujan J. \textbf{69} (2026), article 99.
\url{https://doi.org/10.1007/s11139-026-01372-y};
\url{https://arxiv.org/abs/2409.14502}.
\bibitem{AhlgrenOno} S. Ahlgren and K. Ono,
\emph{Modularity of a certain Calabi--Yau threefold},
Monatsh. Math. \textbf{129} (2000), 177--190.
\url{https://uva.theopenscholar.com/files/ken-ono/files/050.pdf}.
\bibitem{Ahlgren} S. Ahlgren,
\emph{The points of a certain fivefold over finite fields and the twelfth
power of the eta function}, Finite Fields Appl. \textbf{8} (2002), 18--33.
\bibitem{EHMM2} M. Allen, B. Grove, L. Long and F.-T. Tu,
\emph{The explicit hypergeometric-modularity method II},
Res. Math. Sci. \textbf{12} (2025), article~84.
\url{https://doi.org/10.1007/s40687-025-00566-y}.
Preprint: \url{https://arxiv.org/abs/2411.15116}.
\bibitem{DLMF} NIST Digital Library of Mathematical Functions,
Equations 15.10.1, 19.5.1, 19.12.1 and 20.9.2.
\url{https://dlmf.nist.gov/15.10.E1},
\url{https://dlmf.nist.gov/19.5.E1},
\url{https://dlmf.nist.gov/19.12.E1},
\url{https://dlmf.nist.gov/20.9.E2}.

\bibitem{DeligneWII} P. Deligne, \emph{La conjecture de Weil: II},
Publ. Math. IHES \textbf{52} (1980), 137--252,
Corollary~3.3.6 (p.~206); see also Section~3.7.1.
\url{https://publications.ias.edu/sites/default/files/Number40.pdf}.
\bibitem{SchollMotives} A. J. Scholl, \emph{Motives for modular forms},
Invent. Math. \textbf{100} (1990), 419--430.
\url{https://www.dpmms.cam.ac.uk/~ajs1005/preprints/mf.pdf}.
Author's corrections: \url{https://www.dpmms.cam.ac.uk/~ajs1005/preprints/mf_corr.html}.
\bibitem{SchollVanishing} A. J. Scholl,
\emph{Vanishing cycles and non-classical parabolic cohomology},
Invent. Math. \textbf{124} (1996), 503--524.
\url{https://www.dpmms.cam.ac.uk/~ajs1005/preprints/van.pdf}.
Correction: \url{https://www.dpmms.cam.ac.uk/~ajs1005/preprints/van-corr/van-corr.pdf}.
\bibitem{StacksMinimal} The Stacks Project, Lemma 55.10.1, tag 0C6B,
uniqueness of the minimal model.
\url{https://stacks.math.columbia.edu/tag/0C6B}.
\bibitem{AntieauMeier} B. Antieau and L. Meier,
\emph{The Brauer group of the moduli stack of elliptic curves},
Algebra Number Theory \textbf{14} (2020), 2295--2333,
Lemma~4.3 and Proposition~4.5.
\url{https://arxiv.org/pdf/1608.00851}.
\bibitem{WalkeV2} R. M. Walke, \emph{Cubic twists of Legendre moments:
cohomological descent and Frobenius trace formulas}, version 2 (2026).
\url{https://doi.org/10.5281/zenodo.23775155}.
This is the public normalization predecessor of the present research draft.

\bibitem{Ulmer} D. Ulmer, \emph{Conductors of $\ell$-adic representations},
Proc. Amer. Math. Soc. 144 (2016), 2291--2299, \S\S5--9.
\url{https://arxiv.org/abs/1307.4525}.

\bibitem{Nakayama} C. Nakayama, \emph{Nearby cycles for log smooth families},
Compositio Mathematica \textbf{112} (1998), 45--75,
Variant 3.4.1, Proposition 4.2 and Remark 4.2.1.
\url{https://doi.org/10.1023/A:1000327225021}.

\bibitem{SerreRankOne} J.-P. Serre, \emph{Abelian $\ell$-adic Representations and
Elliptic Curves}, 1968; Chapter III, \S1.2, Theorem (Tate), printed p. III-7,
and Appendix. Primary scan:
\url{https://www.math.mcgill.ca/darmon/courses/18-19/gs/serre-mcgill.pdf}.

\bibitem{Yun} Z. Yun, with an appendix by C. Vincent,
\emph{Galois representations attached to moments of Kloosterman sums and
conjectures of Evans}, Compositio Math. \textbf{151} (2015), 68--120,
Sections 2.5.5 and 4.3.1--4.3.4.
\url{https://arxiv.org/abs/1308.3920}.

\bibitem{Haessig} C. D. Haessig,
\emph{Kloosterman sums and Hecke polynomials in characteristics 2 and 3},
Finite Fields Appl. \textbf{76} (2021), 101896, proof of Theorem 2.2.
\url{https://arxiv.org/abs/2012.00570}.

\bibitem{SGA7} P. Deligne and N. Katz (eds.),
\emph{Groupes de monodromie en geometrie algebrique (SGA 7 II)},
Lecture Notes in Math. \textbf{340}, Springer (1973),
Expose XIII, Lemmas 2.1.10--2.1.11.
\url{https://library.slmath.org/nonmsri/sga/sga/pdf/sga7-2.pdf}.

\bibitem{Illusie} L. Illusie, \emph{An overview of the work of K. Fujiwara,
K. Kato, and C. Nakayama on logarithmic etale cohomology}, Asterisque
\textbf{279} (2002), 271--322, Corollaries 8.4.3--8.4.4.
\url{https://www.numdam.org/article/AST_2002__279__271_0.pdf}.
\end{thebibliography}
